\title{DeepSeekMath: Pushing the Limits of Mathematical Reasoning in Open Language Models}

\begin{abstract}
Language models often struggle with mathematical reasoning, largely due to the intricacies and rigor inherent in mathematical tasks. In this work, we present DeepSeekMath~7B, a model built by further pre-training DeepSeek-Coder-Base-v1.5 7B on a corpus comprising 120B math-centric tokens extracted from Common Crawl, alongside datasets including natural language and programming code. DeepSeekMath~7B attains a notable 51.7\% on the challenging MATH benchmark at the competition level, doing so independently of external tools and ensemble methods—approaching the accuracy exhibited by Gemini-Ultra and GPT-4. Furthermore, when applying self-consistency over 64 samples, DeepSeekMath~7B's accuracy rises to 60.9\% on the MATH benchmark. The strengthened mathematical proficiency of DeepSeekMath~stems from two principal innovations: first, the comprehensive utilization of web-scale data via a precisely designed data curation pipeline; second, the development of Group Relative Policy Optimization (GRPO), an adaptation of Proximal Policy Optimization (PPO), which not only sharpens mathematical reasoning but also yields more efficient PPO memory utilization.
\end{abstract}

\section{Introduction}

Large language models (LLMs) have transformed the field of mathematical reasoning within artificial intelligence, leading to remarkable progress on both quantitative reasoning benchmarks \citep{MATH} and geometry reasoning tasks \citep{trinh2024solving}. In addition, these models have demonstrated substantial utility in aiding humans as they tackle challenging mathematical problems \citep{tao}. Nevertheless, state-of-the-art systems like GPT-4 \citep{gpt4} and Gemini-Ultra \citep{gemini} remain proprietary, while existing open-source alternatives exhibit a marked gap in their level of performance.

In this work, we present DeepSeekMath, a specialized language model designed for mathematics that greatly surpasses open-source alternatives in mathematical reasoning and nearly reaches the performance of GPT-4 on standard academic evaluations. Central to our approach is the construction of the DeepSeekMath Corpus, an extensive and high-quality pre-training resource containing 120B math-specific tokens. To build this corpus, we apply a fastText-based classifier \citep{joulin2016fasttext} to filter relevant data from the Common Crawl (CC). Initially, this classifier is trained using positive samples taken from OpenWebMath \citep{openwebmath} and negative samples drawn from a broad spectrum of other websites. We then utilize the trained classifier to identify further positive math instances within the CC, which are subsequently validated and curated by human annotators. Incorporating these refined examples, we retrain the classifier, thereby enhancing its accuracy. Our assessments show that the resulting dataset is of exceptional quality—our foundational model, DeepSeekMath-Base 7B, attains 64.2\% accuracy on GSM8K \citep{gsm8k} and 36.2\% on the challenging MATH benchmark \citep{MATH}, surpassing the Minerva 540B model \citep{minerva}. Notably, because the DeepSeekMath Corpus includes multiple languages, we observe improvements on Chinese mathematical evaluation sets \citep{wei2023cmath,agieval}. We view our methods in mathematical data preparation as a valuable entry point for further research, with considerable potential for future advancements.

DeepSeekMath-Base is built upon DeepSeek-Coder-Base-v1.5 7B \citep{deepseek-coder}, as initializing from a code-focused pretrained model proves more advantageous than employing a standard language model. Additionally, our findings reveal that mathematical training contributes to improved performance not only on mathematical tasks but also across MMLU \citep{mmlu} and BBH benchmarks \citep{bbh}. This suggests that mathematical training strengthens both mathematical proficiency and broader reasoning skills.

Following pre-training, we conduct mathematical instruction tuning on DeepSeekMath-Base utilizing datasets featuring chain-of-thought \citep{cot}, program-of-thought \citep{pot,pal}, and tool-integrated reasoning \citep{tora}. This process yields the DeepSeekMath-Instruct 7B model, which surpasses all other 7B-scale models and demonstrates performance on par with instruction-tuned open-source models at the 70B scale.

In addition, we propose Group Relative Policy Optimization (GRPO), a new reinforcement learning (RL) algorithm inspired by Proximal Policy Optimization (PPO) \citep{schulman2017proximal}. Unlike traditional approaches, GRPO eliminates the need for a critic network and instead leverages group-level scores to estimate the baseline, thereby greatly reducing computational overhead during training. Utilizing only a limited set of English instruction tuning data, GRPO achieves notable gains relative to the strong DeepSeekMath-Instruct model, leading to significant improvements across both in-domain datasets (GSM8K: 82.9\% $\rightarrow$ 88.2\%, MATH: 46.8\% $\rightarrow$ 51.7\%) as well as out-of-domain scenarios (e.g., CMATH: 84.6\% $\rightarrow$ 88.8\%) throughout the RL phase. Moreover, we introduce a unified framework that facilitates the interpretation of several algorithms—including Rejection Sampling Fine-Tuning (RFT) \citep{yuan2023scaling}, Direct Preference Optimization (DPO) \citep{dpo}, PPO, and GRPO—by revealing that all can be represented as variants of either conventional or streamlined RL approaches. We further undertake comprehensive empirical evaluations, contrasting online and offline regimes, outcome-based versus process-based supervision, and single-turn against iterative RL configurations, among other factors, to thoroughly analyze the core components of the unified framework. Finally, we clarify the mechanisms through which our RL methodologies enhance the capabilities of instruction-tuned models, and we outline promising avenues for advancing RL techniques in accordance with this overarching paradigm.

\subsection{Contributions}
We present scalable methods for mathematical pre-training and conduct a thorough investigation and evaluation of reinforcement learning techniques.

\textbf{Math Pre-Training at Scale}

\begin{itemize}[topsep=0pt]
    \item This study demonstrates that the openly available Common Crawl dataset harbors significant mathematical content. Through the use of a rigorously crafted selection pipeline, we were able to create the DeepSeekMath Corpus—a high-quality collection consisting of 120B tokens drawn from web pages with a focus on mathematical material. This corpus is nearly seven times larger than the math web page subset utilized by Minerva \citep{minerva}, and exceeds the size of the recently published OpenWebMath dataset \citep{openwebmath} by a factor of nine.

\item DeepSeekMath-Base 7B, our pre-trained base model, demonstrates mathematical reasoning performance on par with Minerva 540B \citep{minerva}. This suggests that sheer parameter count is not the sole determinant of proficiency in this domain; a relatively compact model, when trained on curated, high-quality datasets, can also deliver robust results.

\item We present the results of our experiments focused on math training. Introducing code training before math training enhances model performance in mathematical problem solving, regardless of whether external tools are employed. This contributes evidence to the ongoing discussion of whether code training benefits reasoning skills, and suggests that, for mathematical reasoning in particular, it does yield improvements.

\item While numerous mathematics-focused works frequently utilize arXiv papers for training, this practice yields no significant enhancements across the various mathematical benchmarks evaluated in this study.
\end{itemize}

\textbf{Exploration and Analysis of Reinforcement Learning}

\begin{itemize}[topsep=0pt]
    \item We present Group Relative Policy Optimization (GRPO), a novel reinforcement learning algorithm that emphasizes efficiency and efficacy. Unlike traditional approaches relying on a critic model, GRPO computes the baseline using group-level scores, thus greatly reducing the computational demands relative to Proximal Policy Optimization (PPO).
    \item Our results indicate that applying GRPO yields substantial improvements to the instruction-tuned DeepSeekMath-Instruct model, solely leveraging the available instruction-tuning dataset. Moreover, we find that this approach promotes gains in out-of-domain generalization during reinforcement learning.
    \item We establish a unified framework to analyze and compare various methods such as RFT, DPO, PPO, and GRPO. Additionally, we undertake a comprehensive series of experiments—including comparisons of online versus offline learning, outcome versus process-level feedback, and single-step versus iterative reinforcement learning—to systematically explore key components of this paradigm.
    \item Building on our unified framework, we investigate underlying mechanisms contributing to reinforcement learning effectiveness, and outline several promising research avenues for further improving reinforcement learning methods in the context of LLMs.
\end{itemize}

\subsection{Summary of Evaluations and Metrics}

\begin{itemize}[topsep=0pt]
    \item \textbf{English and Chinese Mathematical Reasoning}: We perform in-depth evaluations of our models using a range of benchmarks in both English and Chinese, spanning mathematical questions from elementary to collegiate complexity. The English benchmarks utilized are GSM8K \citep{gsm8k}, MATH \citep{MATH}, SAT \citep{llemma}, OCW Courses \citep{minerva}, and MMLU-STEM \citep{mmlu}. For Chinese, assessments are based on MGSM-zh \citep{mgsm}, CMATH \citep{wei2023cmath}, Gaokao-MathCloze \citep{agieval}, and Gaokao-MathQA \citep{agieval}. Model performance is tested both in terms of independently producing complete text-based solutions without external tools, as well as solving tasks through the application of Python.

On evaluations involving English datasets, DeepSeekMath-Base performs on par with the proprietary Minerva 540B model \citep{minerva}, and outperforms all existing open-source base models—including Mistral 7B \citep{mistral} and Llemma-34B \citep{llemma}—by a substantial margin, irrespective of whether these models were math pre-trained. Remarkably, DeepSeekMath-Base achieves even better results on Chinese-language benchmarks. This advantage is likely attributable to our approach of not restricting the math pre-training corpus to English—as previous studies have done \citep{minerva,llemma}—but instead incorporating high-quality data from multiple languages. Building on this, mathematical instruction tuning and reinforcement learning yield DeepSeekMath-Instruct and DeepSeekMath-RL, which together deliver state-of-the-art results, with accuracy surpassing 50\% on the challenging MATH competition dataset—marking a first for open-source models.

\item \textbf{Formal Mathematics}: DeepSeekMath-Base is assessed through the informal-to-formal theorem proving task as outlined in \citep{dsp_proof}, specifically employing miniF2F \citep{minif2f} with Isabelle \citep{isabelle} as the designated proof assistant. The results indicate that DeepSeekMath-Base achieves notable performance in few-shot autoformalization scenarios.
\item \textbf{Natural Language Understanding, Reasoning, and Code}: To thoroughly assess the models' proficiency in general comprehension, logical reasoning, and programming, we measure DeepSeekMath-Base on the Massive Multitask Language Understanding (MMLU) benchmark \citep{mmlu}, which features 57 multiple-choice problems spanning a range of topics; the BIG-Bench Hard (BBH) suite \citep{bbh}, which includes 23 complex problems requiring advanced reasoning; as well as HumanEval \citep{codex} and MBPP \citep{mbpp}, both commonly adopted for evaluating code generation models. The results suggest that mathematical pre-training substantially enhances both language understanding and reasoning abilities.

\section{Math Pre-Training}

\subsection{Data Collection and Decontamination}  
This section describes the methodology for assembling the DeepSeekMath Corpus utilizing data from Common Crawl. As illustrated in Figure \ref{fig:data_collect}, we introduce an iterative pipeline designed to efficiently amass a substantial mathematics-focused corpus, beginning with a seed set—namely, a limited yet high-caliber collection of math-oriented datasets. It is important to emphasize that this method is generalizable and can be extended to other fields, such as coding.

Initially, we select OpenWebMath \citep{openwebmath}, a repository comprised of well-curated mathematical texts from the web, to serve as our foundational seed corpus. Leveraging this resource, we proceed to train a fastText model \citep{joulin2016fasttext} for the retrieval of additional mathematical web documents resembling those found in OpenWebMath. For this purpose, we randomly draw 500,000 items from the seed corpus to serve as positive examples, while another 500,000 web pages are randomly taken from Common Crawl to serve as negative examples. Utilizing an open-source toolkit\footnote{\url{https://fasttext.cc}} for fastText training, we specify a vector size of 256, set the learning rate to 0.1, establish a maximum word n-gram length of 3, require at least 3 occurrences per word, and run for 3 epochs. To narrow the scope of Common Crawl, we apply both URL-based deduplication and near-deduplication methodologies, yielding a set of 40 billion unique HTML documents. Employing the fastText model, we search for mathematical pages within this deduplicated corpus. To mitigate the inclusion of low-quality mathematical text, the recovered documents are ordered by the scores assigned by the model, and only those ranking highest are retained. We then determine the suitable data retention amount by evaluating pre-training performance using subsets of the top 40B, 80B, 120B, and 160B tokens. For the initial pass, we opt to retain the top 40B tokens.

Following the initial round of data collection, a significant number of mathematical web pages are left uncollected. This shortcoming is largely due to the limited diversity present in the positive example set used to train the fastText model. To address this, we seek out further mathematical web resources to diversify and expand the seed corpus, aiming to enhance fastText model performance. Specifically, the entire Common Crawl is partitioned into separate domains, where each domain consists of web pages sharing the same base URL. For every domain, we determine the proportion of pages acquired during the first pass. Domains in which more than 10\% of pages are collected are labeled as math-related (for instance, \url{mathoverflow.net}).
Next, we perform manual annotation of URLs representing mathematical content within these math-oriented domains (such as \url{mathoverflow.net/questions}). Any remaining web pages that are connected to these URLs, but were not previously collected, are incorporated into the seed corpus. With this methodology, we are able to gather a wider variety of positive samples, subsequently enabling the training of a more robust fastText model that can identify a greater quantity of mathematical data in subsequent rounds. After repeating this collection process over four cycles, we obtain a total of 35.5M mathematical web pages, amounting to 120B tokens. During the fourth round, we observe that almost 98\% of the mathematical content had already been retrieved by the third round, which leads us to conclude the data collection process.

In order to prevent contamination of our benchmarks, we adopt the methodology outlined in \cite{deepseek-coder} to exclude web pages that include content from English mathematics benchmarks, such as GSM8K \citep{gsm8k} and MATH \citep{MATH}, as well as Chinese benchmarks like CMATH \citep{wei2023cmath} and AGIEval \citep{agieval}. Specifically, we apply the following filtering procedure: any text fragment containing a sequence of 10 contiguous words that is an exact match to a substring present in these benchmark datasets is discarded from our math training data. For benchmark passages containing fewer than 10-gram sequences but at least 3-gram sequences, we carry out the same procedure by applying exact matching with the available n-grams, removing any web pages that contain these matches.

\subsection{Validating the Quality of the DeepSeekMath~Corpus}

We run pre-training experiments to investigate how the DeepSeekMath~Corpus is compared with the recently released math-training corpora:

\begin{itemize}[topsep=0pt]
    \item \textbf{MathPile} \citep{mathpile}: This dataset is a collection compiled from various resources, including textbooks, Wikipedia, ProofWiki, CommonCrawl, StackExchange, and arXiv, and comprises a total of 8.9B tokens; notably, over 85\% of the material originates from arXiv.
    \item \textbf{OpenWebMath} \citep{openwebmath}: Composed of CommonCrawl data that has been specifically filtered to retain mathematical content, this dataset amounts to 13.6B tokens.
    \item \textbf{Proof-Pile-2} \citep{llemma}: Encompassing OpenWebMath, AlgebraicStack (which consists of 10.3B tokens from mathematical code), and arXiv documents (with 28.0B tokens), this corpus is dedicated to mathematical text. For experiments involving Proof-Pile-2, we adhere to the protocol from \cite{llemma}, applying an arXiv:Web:Code composition ratio of 2:4:1.
\end{itemize}

\subsubsection{Training Setting}
\label{sec:quality-policy}
For our experiments, we fine-tune a general-purpose pre-trained language model comprising 1.3 billion parameters, structured analogously to the DeepSeek LLMs \citep{deepseek-llm}, which we refer to as DeepSeek-LLM 1.3B. Each individual mathematical dataset is used to separately train the model for a total of 150 billion tokens. All model training is performed using the streamlined and resource-efficient HAI-LLM framework \citep{haillm}. We adhere to the established training protocol from DeepSeek LLMs, employing the AdamW optimizer \citep{adamW} configured with $\beta_1=0.9$, $\beta_2=0.95$, and a $\mathrm{weight\_decay}$ of 0.1. The learning rate is scheduled in multiple phases: it is ramped up to its maximum after 2,000 warmup steps, then reduced to 31.6\% of its peak at 80\% completion of training, and finally decayed to 10.0\% of the peak at 90\% training progress. The upper bound for the learning rate is set at 5.3e-4. Additionally, the batch size is configured to encompass 4 million tokens and a context window of 4K tokens is maintained.

\subsubsection{Evaluation Results}

\textbf{The DeepSeekMath~Corpus is of high quality, covers multilingual mathematical content, and is the largest in size.}

\begin{itemize}[topsep=0pt]
    \item \textbf{High-quality}:
    We assess downstream effectiveness across 8 mathematical tasks employing few-shot chain-of-thought prompting \cite{cot}. Table \ref{tab:corpora_comparison} illustrates that models trained on the DeepSeekMath~Corpus consistently outperform alternatives. Moreover, as depicted in Figure \ref{fig:corpora_comparisons}, the model trained on the DeepSeekMath Corpus surpasses the Proof-Pile-2-trained counterpart after 50B tokens (equivalent to one full epoch on Proof-Pile-2), highlighting the overall superior quality of data within the DeepSeekMath Corpus.
    \item \textbf{Multilingual}:
    The DeepSeekMath~Corpus contains data spanning multiple languages, with English and Chinese constituting the largest proportions. Evidence from Table \ref{tab:corpora_comparison} demonstrates that models utilizing the DeepSeekMath~Corpus display enhanced mathematical reasoning abilities in both English and Chinese. In contrast, models built with predominantly English-oriented mathematical corpora experience minimal, or even negative, gains when evaluated on Chinese mathematical reasoning tasks.
    \item \textbf{Large-scale}:
    DeepSeekMath~Corpus significantly exceeds the scale of other existing mathematical datasets. As shown in Figure \ref{fig:corpora_comparisons}, DeepSeek-LLM 1.3B, after being trained on DeepSeekMath~Corpus, exhibits a more accelerated and sustained learning curve. This stands in stark contrast to smaller, baseline corpora that have been cycled through several times during training and whose performance improvements quickly saturate.
\end{itemize}

\subsection{Training and Evaluating DeepSeekMath-Base 7B}

This section presents DeepSeekMath-Base 7B, a foundational model designed with notable reasoning capabilities, particularly within the field of mathematics. The model’s initialization leverages DeepSeek-Coder-Base-v1.5 7B \citep{deepseek-coder}, and its training spans 500B tokens. The composition of training data consists of 56\% DeepSeekMath Corpus, 4\% AlgebraicStack, 10\% arXiv, 20\% code from Github, and 10\% natural language samples drawn from Common Crawl, encompassing both English and Chinese texts. Training procedures largely adhere to the configuration detailed in Section \ref{sec:quality-policy}, with the primary distinctions being a learning rate capped at 4.2e-4 and a batch size set to 10M tokens.

We present an in-depth evaluation of DeepSeekMath-Base 7B’s mathematical proficiency, emphasizing its capacity to independently generate complete mathematical solutions, utilize tools to address mathematical questions, and execute formal theorem proving. In addition to its mathematical aptitudes, we further examine the base model’s broader competencies, encompassing natural language understanding, logical reasoning, and coding abilities.

\paragraph{Mathematical Problem Solving with Step-by-Step Reasoning}

We assess the ability of DeepSeekMath-Base to solve mathematical problems through few-shot chain-of-thought prompting \citep{cot}, utilizing eight benchmarks in both English and Chinese. These benchmarks span a range of tasks, including quantitative reasoning (such as GSM8K \citep{gsm8k}, MATH \citep{MATH}, and CMATH \citep{wei2023cmath}) and multiple-choice formats (for example, MMLU-STEM \citep{mmlu} and Gaokao-MathQA \citep{agieval}). Collectively, these benchmarks represent a wide array of mathematical domains, with difficulty ranging from elementary school up to college-level mathematics.

Table \ref{tab:base_math_cot} demonstrates that DeepSeekMath-Base 7B achieves the highest scores across all eight evaluated benchmarks when compared to other open-source base models, such as the popular Mistral 7B \citep{mistral} and the newer Llemma 34B \citep{llemma}, which received additional training on Proof-Pile-2 \citep{llemma} for mathematical capabilities. Significantly, DeepSeekMath-Base exceeds the top-performing open-source base models by more than 10\% absolute on the rigorous, competition-level MATH dataset. Furthermore, it surpasses the closed-source Minerva 540B \citep{minerva}, which is built upon PaLM \citep{palm} and fine-tuned using mathematical literature, despite Minerva’s considerably greater parameter count—being 77 times larger.

\paragraph{Mathematical Problem Solving with Tool Use} To assess program-supported mathematical reasoning, we conduct experiments on GSM8K and MATH datasets employing few-shot program-of-thought prompting \citep{pot,pal}. In this approach, models generate solutions by composing Python scripts that may leverage external libraries like \textit{math} and \textit{sympy} for complex numerical tasks. The output produced by executing these scripts is regarded as the final answer. According to Table \ref{tab:base_math_pot_proof}, DeepSeekMath-Base 7B surpasses the previous best model, Llemma 34B, in performance.

\paragraph{Formal Mathematics}

The automation of formal proofs offers significant advantages by improving the correctness and dependability of mathematical proofs, as well as boosting productivity—a topic that has received growing interest in recent years. In this context, we assess DeepSeekMath-Base 7B on the informal-to-formal proving task presented in \citep{dsp_proof}, which involves producing a formal proof given an informal problem description, its formalized version, and an informal argument. Our experiments utilize miniF2F \citep{minif2f}, a standard benchmark focusing on formal mathematics at the Olympiad level. For each problem, we employ few-shot prompting to prompt Isabelle to generate a corresponding formal proof. In line with \cite{dsp_proof}, we prompt language models to output proof outlines, and subsequently apply the automated prover Sledgehammer \citep{sledgehammer} to complete any unfinished portions of the argument. According to the results displayed in Table \ref{tab:base_math_pot_proof}, DeepSeekMath-Base 7B achieves impressive results in the automatic formalization of proofs.

\paragraph{Natural Language Understanding, Reasoning, and Code}

We assess the capabilities of the model in natural language understanding using MMLU \citep{mmlu}, in reasoning tasks through BBH \citep{bbh}, and in programming proficiency via HumanEval \citep{codex} and MBPP \citep{mbpp}. According to the results presented in Table \ref{tab:understanding_reasoning_code}, DeepSeekMath-Base 7B achieves notable improvements on both the MMLU and BBH datasets compared to its predecessor, DeepSeek-Coder-Base-v1.5 \citep{deepseek-coder}, thereby demonstrating the beneficial effects of specialized math training on linguistic comprehension and reasoning skills. Furthermore, by incorporating code-related data during continued pretraining, DeepSeekMath-Base 7B successfully preserves the coding performance of DeepSeek-Coder-Base-v1.5 across the evaluated programming tasks. In summary, DeepSeekMath-Base 7B delivers superior outcomes relative to the general-purpose Mistral 7B model \citep{mistral} across all three reasoning and coding evaluations.

\section{Supervised Fine-Tuning}

\subsection{SFT Data Curation}

A mathematical instruction-tuning dataset is developed, encompassing a diverse range of English and Chinese mathematical problems that span multiple domains and present different levels of difficulty. Each problem is accompanied by a solution formatted as chain-of-thought (CoT) \citep{cot}, program-of-thought (PoT) \citep{pot,pal}, or tool-integrated reasoning \citep{tora}. In total, the dataset comprises 776K training instances.

\begin{itemize}[topsep=0pt]
    \item \textbf{English mathematical datasets}: We provide tool-assisted solution annotations for problems sourced from GSM8K and MATH, and further include a selected subset from MathInstruct \citep{MathInstruct} in addition to the Lila-OOD training set \citep{lila}, both of which feature solutions constructed using CoT or PoT approaches. The English dataset spans a broad range of mathematical domains, including but not limited to algebra, probability, number theory, calculus, and geometry.
    \item \textbf{Chinese mathematical datasets}: Our collection consists of K-12 mathematics problems in Chinese, encompassing 76 different sub-topics such as linear equations; each problem is annotated with both CoT and tool-integrated reasoning solutions.
\end{itemize}

\subsection{Training and Evaluating DeepSeekMath-Instruct 7B}

This section presents DeepSeekMath-Instruct 7B, a model refined with mathematical instruction tuning on top of DeepSeekMath-Base. To prepare the training data, samples are combined in random order, creating input sequences up to a maximum context length of 4K tokens. The model is trained for 500 steps using a batch size of 256, maintaining a fixed learning rate of 5e-5 throughout the process.

We evaluate models' mathematical performance both without and with tool use, on 4 quantitative reasoning benchmarks in English and Chinese. We benchmark our model against the leading models of the time:

\begin{itemize}[topsep=0pt]
    \item \textbf{Closed-source models} encompass: (1) the GPT family, with the most advanced models being GPT-4 \citep{gpt4} and GPT-4 Code Interpreter~\footnote{\url{https://openai.com/blog/chatgpt-plugins\#\#code-interpreter}}, (2) Gemini Ultra and Pro \citep{gemini}, (3) Inflection-2 \citep{inflection-2}, (4) Grok-1~\footnote{\url{https://x.ai/model-card}},
    as well as some recent offerings from Chinese developers, including (5) Baichuan-3~\footnote{\url{https://www.baichuan-ai.com}},
    and (6) the newest version of GLM-4~\footnote{\url{https://open.bigmodel.cn/dev/api\#glm-4}}

    from the GLM family \citep{glm}.

These models serve a broad range of applications, and most have been subjected to multiple alignment processes.  
\item \textbf{Open-source models} span: standard models such as (1) DeepSeek-LLM-Chat 67B \citep{deepseek-llm}, (2) Qwen 72B \citep{qwen}, (3) SeaLLM-v2 7B \citep{seallm}, and (4) ChatGLM3 6B \citep{chatglm3}; in addition, there are models focused on mathematical tasks, such as (5) InternLM2-Math 20B~\footnote{\url{https://github.com/InternLM/InternLM-Math}}, an extension of InternLM2 with specialized math-oriented pre-training and subsequent instruction tuning, (6) Math-Shepherd-Mistral 7B, which enhances Mistral 7B \citep{mistral} with PPO-based training \citep{schulman2017proximal} guided by a process-supervised reward model, (7) the WizardMath family \citep{wizardmath} which advances mathematical reasoning capabilities in Mistral 7B and Llama-2 70B \citep{llama2} via evolve-instruct (a variant of instruction tuning employing instructions created through AI evolution) and PPO, using primarily problems from GSM8K and MATH during training, (8) MetaMath 70B \citep{metamath}, produced by fine-tuning Llama-2 70B on an enriched collection combining GSM8K and MATH, (9) ToRA 34B \cite{tora}, in which CodeLlama 34B is adapted for tool-augmented mathematical reasoning, and (10) MAmmoTH 70B \citep{MathInstruct}, where Llama-2 70B is instruction-tuned utilizing MathInstruct data.

Table \ref{tab:sft_rl_math} illustrates that, when evaluated in scenarios where tool use is prohibited, DeepSeekMath-Instruct 7B exhibits robust capabilities in stepwise reasoning. Importantly, on the competition-oriented MATH benchmark, our model outperforms every open-source alternative as well as most closed-source systems (such as Inflection-2 and Gemini Pro) by a margin of at least 9 percentage points. This advantage holds even in comparison to models with significantly larger architectures (e.g., Qwen 72B) or those specifically refined with math-centric reinforcement learning approaches (such as WizardMath-v1.1 7B). Although DeepSeekMath-Instruct is on par with Chinese proprietary systems like GLM-4 and Baichuan-3 on the MATH dataset, it still falls short of the performance reached by GPT-4 and Gemini Ultra.

In an evaluation scenario permitting models to combine natural language reasoning with programmatic tool usage for solving tasks, DeepSeekMath-Instruct 7B attains nearly 60\% accuracy on MATH, outperforming all previously available open-source models. Across additional benchmarks, our model demonstrates performance on par with DeepSeek-LLM-Chat 67B, the earlier state-of-the-art model that boasts a parameter count ten times greater.  
\section{Reinforcement Learning}

\subsection{Group Relative Policy Optimization} Previous research demonstrates that reinforcement learning (RL) is a valuable method for enhancing the mathematical reasoning capabilities of LLMs subsequent to the Supervised Fine-Tuning (SFT) phase \citep{wang2023math,wizardmath}. Here, we present our novel and practical RL approach, Group Relative Policy Optimization (GRPO).

\subsubsection{From PPO to GRPO} Proximal Policy Optimization (PPO) \citep{schulman2017proximal} is an actor-critic RL algorithm that is widely used in the RL fine-tuning stage of LLMs \citep{ouyang2022training}. In particular, it optimizes LLMs by maximizing the following surrogate objective:

\begin{equation}
\footnotesize
    \mathcal{J}_{PPO}(\theta) = \mathbb{E}{[q \sim P(Q), o \sim \pi_{\theta_{old}}(O|q)]} \frac{1}{|o|} \sum_{t=1}^{|o|} \min \left[ \frac{\pi_\theta(o_{t} | q, o_{<t})}{\pi_{\theta_{old}}(o_{t} | q, o_{<t})} A_{t}, \text{clip} \left( \frac{\pi_\theta(o_{t} | q, o_{<t})}{\pi_{\theta_{old}}(o_{t} | q, o_{<t})}, 1 - \varepsilon, 1 + \varepsilon \right)  A_{t} \right] ,
\end{equation}

The objective function for \textit{Proximal Policy Optimization (PPO)} is constructed by taking the expectation over queries $q$ sampled from $P(Q)$ and outputs $o$ generated by the previous policy $\pi_{\theta_{old}}(O|q)$. This objective incorporates a per-token minimum of the importance-weighted advantage and a version of this ratio that is clipped between $1 - \varepsilon$ and $1 + \varepsilon$ to limit the update size. Specifically, for each token position $t$, the function averages this minimum across all tokens in $o$, ensuring stable policy improvement through conservative likelihood ratio updates.

In this context, $\pi_{\theta}$ represents the present policy model, while $\pi_{\theta_{old}}$ denotes the preceding version. The variables $q$ and $o$ correspond to questions and model outputs, which are drawn from the question dataset and generated by the previous policy $\pi_{\theta_{old}}$, respectively. The term $\varepsilon$ acts as a hyper-parameter in PPO associated with clipping, which serves to enhance the stability of the training process. The advantage term $A_t$ is estimated through the use of Generalized Advantage Estimation (GAE) \citep{gae}, calculated from the sequence of rewards $\{r_{\ge t}\}$ in conjunction with the value predicted by a learned function $V_{\psi}$. Consequently, PPO necessitates concurrent training of the value function and the policy model. To prevent excessive optimization of the reward function, it is standard to incorporate a token-level KL penalty into the reward at every step, relative to a reference model, as detailed in \citep{ouyang2022training}; that is,

\begin{equation}
    r_{t} = r_\varphi(q, o_{\le t}) - \beta \log\frac{\pi_{\theta}(o_{t}|q, o_{<t})}{\pi_{ref}(o_{t}|q, o_{<t})},
\label{eq:PPO-reward}
\end{equation}

Here, $r_\varphi$ denotes the reward model, $\pi_{ref}$ corresponds to the reference model—typically the original SFT model—and $\beta$ represents the KL penalty coefficient.

As the value function employed in PPO is typically another model of comparable size as the policy model, it brings a substantial memory and computational burden. Additionally, during RL training, the value function is treated as a baseline in the calculation of the advantage for variance reduction. While in the LLM context, usually only the last token is assigned a reward score by the reward model, which may complicate the training of a value function that is accurate at each token. To address this, as shown in Figure \ref{fig:grpo}, we propose Group Relative Policy Optimization (GRPO), which obviates the need for additional value function approximation as in PPO, and instead uses the average reward of multiple sampled outputs, produced in response to the same question, as the baseline. More specifically, for each question $q$, GRPO samples a group of outputs $\{o_1, o_2, \cdots, o_G\}$  from the old policy  $\pi_{\theta_{old}}$  and then optimizes the policy model by maximizing the following objective:

\begin{equation}
\footnotesize
\begin{split}
    \mathcal{J}_{GRPO}(\theta) &= \mathbb{E}{[q \sim P(Q), \{o_i\}_{i=1}^G \sim \pi_{\theta_{old}}(O|q)]} \\
    & \frac{1}{G} \sum_{i=1}^G \frac{1}{|o_i|} \sum_{t=1}^{|o_i|} \bigg\{ \min \Big[ \frac{\pi_\theta(o_{i,t} | q, o_{i,<t})}{\pi_{\theta_{old}}(o_{i,t} | q, o_{i,<t})} \hat{A}_{i,t}, \\
    & \qquad\qquad \text{clip} \Big( \frac{\pi_\theta(o_{i,t} | q, o_{i,<t})}{\pi_{\theta_{old}}(o_{i,t} | q, o_{i,<t})}, 1 - \varepsilon, 1 + \varepsilon \Big) \hat{A}_{i,t} \Big] \\
    & \qquad\qquad - \beta \mathbb{D}_{KL}\left[\pi_{\theta} || \pi_{ref}\right] \bigg\} ,
\end{split}
\label{eq:GRPO-obj}
\end{equation}

where $\varepsilon$ and $\beta$ are hyper-parameters, and $\hat{A}_{i,t}$  is the advantage calculated based on relative rewards of the outputs inside each group only, which will be detailed in the following subsections. The group relative way that GRPO leverages to calculate the advantages, aligns well with the comparative nature of rewards models,  as reward models are typically trained on datasets of comparisons between outputs on the same question. Also note that, instead of adding KL penalty in the reward, GRPO regularizes by directly adding the KL divergence between the trained policy and the reference policy to the loss, avoiding complicating the calculation of  $\hat{A}_{i,t}$. And different from the KL penalty term used in (\ref{eq:PPO-reward}), we estimate the KL divergence with the following unbiased estimator \citep{kl_approx}:

\begin{equation}
\small
    \mathbb{D}_{KL}\left[\pi_{\theta} || \pi_{ref}\right] = \frac{\pi_{ref}(o_{i,t}|q,o_{i,<t})}{\pi_{\theta}(o_{i,t}|q,o_{i,<t})}- \log\frac{\pi_{ref}(o_{i,t}|q,o_{i,<t})}{\pi_{\theta}(o_{i,t}|q,o_{i,<t})} - 1,
\end{equation}

which is guaranteed to be positive.

\begin{algorithm}[t]
  \small
  \caption{Iterative Group Relative Policy Optimization}
  \textbf{Input} initial policy network $\pi_{\theta_{\text{init}}}$; reward evaluators $r_\varphi$; collection of task prompts $\mathcal{D}$;
  hyperparameters $\varepsilon$, $\beta$, $\mu$
  \begin{algorithmic}[1]
    \State Initialize policy $\pi_\theta \leftarrow \pi_{\theta_{\text{init}}}$
    \For{iteration = 1, \dots, I}
       \State Set reference model $\pi_{ref} \leftarrow \pi_{\theta}$
      \For{step = 1, \dots, M}
      \State Draw mini-batch $\mathcal{D}_b$ from $\mathcal{D}$
      \State Copy current policy to $\pi_{\theta_{old}} \leftarrow \pi_{\theta}$
      \State For every question $q$ in $\mathcal{D}_b$, sample $G$ responses $\{o_i\}_{i=1}^G \sim \pi_{\theta_{old}} (\cdot \mid q) $
      \State For each generated response $o_i$, calculate reward $r_i$ with $r_\varphi$
      \State For each sampled output, estimate $\hat{A}_{i,t}$ at token $t$ via group-based relative advantage calculation.
      \For{GRPO iteration = 1, \dots, $\mu$}
        \State Adjust policy $\pi_\theta$ by optimizing the GRPO loss (Equation \ref{eq:GC-GRPO})
      \EndFor
    \EndFor
    \State Periodically refine $r_\varphi$ via replay buffer updates and continued training.
    \EndFor
  \end{algorithmic}
  \textbf{Output} $\pi_\theta$
  \label{alg:iter-grpo}
\end{algorithm}

\subsubsection{Outcome Supervision RL with GRPO} Specifically, given a question $q$, the old policy $\pi_{\theta_{old}}$ generates a collection of outputs $\{o_1, o_2, \cdots, o_G\}$. Each output is evaluated using a reward model, resulting in a set of rewards $\mathbf{r} = \{r_1, r_2, \cdots, r_G\}$. These rewards are then standardized by removing the mean and scaling by the standard deviation within the group. Under outcome supervision, the resulting normalized reward for each output $o_i$ is assigned uniformly to all its tokens as the advantage, namely $\hat{A}_{i, t} = \widetilde{r}_i = \frac{r_i- {\rm mean}(\mathbf{r})}{{\rm std}(\mathbf{r})}$. The policy is updated to maximize the objective function presented in equation (\ref{eq:GRPO-obj}).

\subsubsection{Process Supervision RL with GRPO} Outcome supervision limits reward feedback to only the final output, which can be inadequate for effectively guiding the policy, especially in intricate mathematical problem-solving. Building on the approach of \cite{wang2023math}, we adopt process supervision, where feedback is provided after each reasoning stage. More specifically, given a question $q$ and a set of $G$ sampled sequences $\{o_1, o_2, \cdots, o_G\}$, we employ a process reward model that assigns a reward to every step within each generated sequence, resulting in the collection of rewards: $\mathbf{R} = \{ \{r_1^{index(1)},\cdots,r_1^{index(K_1)}\}, \cdots,  \{r_G^{index(1)},\cdots,r_G^{index(K_G)}\} \}$, where $index(j)$ denotes the position of the end token for the $j$-th reasoning step, and $K_i$ represents the total number of reasoning steps for the $i$-th output. To standardize the reward signals, we apply normalization using the overall mean and standard deviation across $\mathbf{R}$, calculated as $\widetilde{r}_i^{index(j)} = \frac{r_i^{index(j)} - {\rm mean(\mathbf{R})}}{{\rm std(\mathbf{R})}}$.

Process supervision then determines the advantage for each token as the cumulative sum of normalized rewards for all subsequent steps, formally given by $\hat{A}_{i, t} = \sum_{index(j) \ge t} \widetilde{r}_i^{index(j)}$. The policy is subsequently updated by maximizing the objective function introduced in equation (\ref{eq:GRPO-obj}).

\subsubsection{Iterative RL with GRPO} During reinforcement learning, the previously used reward model may become inadequate for guiding the evolving policy model. To address this, we investigate the iterative RL framework utilizing GRPO. As detailed in Algorithm \ref{alg:iter-grpo}, iterative GRPO entails creating fresh datasets for the reward model by sampling from the current policy model. The reward model is updated continually through a replay mechanism that retains 10\% of prior data. Subsequently, the policy model is designated as the reference, and it is further optimized using feedback from the newly updated reward model.

\subsection{Training and Evaluating DeepSeekMath-RL}

Our reinforcement learning (RL) experiments utilize DeepSeekMath-Instruct 7B as the foundation. For RL training, we use chain-of-thought questions from GSM8K and MATH present within the SFT dataset, amounting to approximately 144,000 items. To specifically evaluate the effect of RL on benchmarks that have no exposure during the RL procedure, we deliberately exclude all other SFT questions. The procedure for assembling the reward model training set is consistent with that of \citep{wang2023math}. The reward model is initialized from DeepSeekMath-Base 7B and optimized using a learning rate of $2\times10^{-5}$. During GRPO training, the learning rate for the policy model is fixed at $1\times10^{-6}$, with a KL divergence coefficient of 0.04. We generate $64$ samples per question, restrict sequences to a maximum length of 1024 tokens, and set a training batch size of 1024. Each iteration, the policy model undergoes one update following exploration. For evaluation, we apply the same benchmark protocols as used for DeepSeekMath-Instruct 7B. In this setup, the tasks involving GSM8K and MATH with chain-of-thought responses are treated as in-domain, while all other benchmarks are considered out-of-domain.

Table \ref{tab:sft_rl_math} presents comparative results for both open-source and proprietary models employing chain-of-thought and tool-integrated reasoning across English and Chinese evaluation suites. The key observations are as follows: 1) DeepSeekMath-RL 7B achieves accuracy rates of 88.2\% on GSM8K and 51.7\% on MATH when applying chain-of-thought reasoning, outperforming every other open-source model in the 7B–70B parameter range and exceeding the results of most closed-source competitors. 2) Notably, the training of DeepSeekMath-RL 7B was restricted to chain-of-thought-format instruction tuning datasets from GSM8K and MATH, with DeepSeekMath-Instruct 7B as its initialization point. In spite of its limited training domain, it demonstrates superior performance to DeepSeekMath-Instruct 7B across all benchmarked criteria, thereby highlighting the substantial benefits conferred by reinforcement learning.

\section{Discussion}
Here, we present the insights gained from our investigations in both pre-training and reinforcement learning experiments.

\subsection{Lessons Learnt in Pre-Training}

To begin, we recount our observations from the pre-training phase. Except where noted, we follow the training configurations described in Section \ref{sec:quality-policy}. Notably, in this section, references to the DeepSeekMath Corpus pertain to an 89B-token dataset compiled during the second round of data collection.

\subsubsection{Code Training Benefits Mathematical Reasoning}

An often-cited but unproven conjecture posits that training on code enhances reasoning capabilities. In this work, we seek to address this claim in part, focusing on mathematics: specifically, we demonstrate that code training augments models’ mathematical reasoning skills, irrespective of whether external tools are utilized.

To study how code training affects mathematical reasoning, we experimented with the following two-stage training and one-stage training settings:

\textbf{Two-Stage Training}

\begin{itemize}[topsep=0pt]
    \item \textbf{Code Training for 400B Tokens $\rightarrow$ Math Training for 150B Tokens}: Our approach involves pretraining DeepSeek-LLM 1.3B on 400B tokens from code data, which is then succeeded by a stage of 150B math tokens for continued training.
    \item \textbf{General Training for 400B Tokens $\rightarrow$ Math Training for 150B Tokens}: For comparison, we run a control experiment where the initial pretraining is conducted on 400B general tokens drawn from a broad corpus compiled by DeepSeek-AI, as opposed to code tokens, to assess whether initial code token exposure provides an edge in developing mathematical reasoning capabilities.
\end{itemize}

\textbf{One-Stage Training}

\begin{itemize}[topsep=0pt]
    \item \textbf{150B Tokens of Mathematical Pretraining}: DeepSeek-LLM 1.3B is pretrained on a dataset comprising 150 billion mathematics tokens.
    \item \textbf{Combined Pretraining with 400B Code Tokens and 150B Math Tokens}: Previous attempts to introduce math-focused training after code pretraining have led to reduced code-related performance. Here, we explore whether integrating math and code data in a single-phase, mixed-token training approach can both enhance mathematical abilities and counteract the effect of catastrophic forgetting in code tasks.
\end{itemize}

\paragraph{Results}
The downstream outcomes for various training configurations are presented in Table \ref{tab:code-math} and Table \ref{tab:code-math-general-eval}.

Incorporating code training supports program-aided mathematical reasoning in both two-stage and one-stage training scenarios. As indicated in Table \ref{tab:code-math}, within the two-stage approach, engaging in code-focused training on its own notably improves performance on GSM8K and MATH tasks leveraging Python. When additional mathematical training is performed in the second phase, even greater gains are observed. Notably, the one-stage paradigm, which involves a combined presentation of code and math tokens, not only prevents the catastrophic forgetting often encountered with two-stage training but also enhances proficiency in coding tasks (see Table \ref{tab:code-math-general-eval}) as well as program-aided mathematical reasoning (refer to Table \ref{tab:code-math}).

Training on code also enhances mathematical reasoning capabilities even in the absence of tool assistance. In the two-stage training framework, introducing code training as the first phase yields notable, albeit moderate, gains. This phase additionally accelerates progress in later math training, culminating in the highest overall performance. In contrast, jointly training with code and math tokens in a single phase appears to diminish mathematical reasoning in scenarios where tools are not employed. One possible explanation is that DeepSeek-LLM 1.3B’s restricted parameter size may prevent it from adequately integrating both code and mathematical content at once.

\subsubsection{ArXiv Papers Seem Ineffective in Improving Mathematical Reasoning}

ArXiv papers are commonly included as a component of math pre-training data \citep{minerva,gpt-f,llemma,mathpile}. However, detailed analysis regarding their impact on mathematical reasoning has not been extensively conducted. Perhaps counter-intuitively, according to our experiments, arXiv papers seem ineffective in improving mathematical reasoning. We experiment with models of different sizes, including DeepSeek-LLM 1.3B and DeepSeek-Coder-Base-v1.5 7B \citep{deepseek-coder}, using arXiv corpora that underwent varied processing pipelines:

\begin{itemize}[topsep=0pt]
    \item \textbf{MathPile} \citep{mathpile}: an 8.9-billion-token dataset assembled by applying various cleaning and filtering heuristics, with more than 85\% of its content originating from scientific papers on arXiv;
    \item \textbf{ArXiv-RedPajama} \citep{redpajama}: a collection comprising all arXiv LaTeX documents with sections such as preambles, comments, macros, and bibliographies excised, amounting to a total of 28.0 billion tokens.
\end{itemize}

In our study, we independently train DeepSeek-LLM 1.3B on 150B tokens and DeepSeek-Coder-Base-v1.5 7B on 40B tokens using each arXiv-based dataset. Our findings indicate that content from arXiv papers does not contribute effectively to enhancing mathematical reasoning skills. When exposed exclusively to arXiv text, both models fail to achieve significant gains and in some cases show declining performance across a range of mathematical evaluation benchmarks of varying levels of difficulty used here. These assessments include quantitative reasoning tasks such as GSM8K and MATH (see Table \ref{tab:arxiv-cot}), multiple-choice assessments like MMLU-STEM (Table \ref{tab:arxiv-cot}), and benchmarks of formal mathematics including miniF2F (see Table \ref{tab:arxiv_atp}).

However, this conclusion has its limitations and should be taken with a grain of salt. We have not yet studied:

\begin{itemize}[topsep=0pt]
    \item The influence of arXiv tokens on certain math-oriented tasks not examined in this study, for example, the informalization of theorems, where formal statements or proofs are transformed into their informal counterparts;
    \item The potential outcomes of integrating arXiv tokens with additional categories of data;
    \item If the advantages derived from arXiv papers persist or become more pronounced as model sizes increase.
\end{itemize}

Thus, further exploration is required, which we leave for future studies.

\subsection{Insights of Reinforcement Learning}
\subsubsection{Towards to a Unified Paradigm} In this part, we introduce a cohesive framework to systematically examine various training approaches, including SFT, RFT, DPO, PPO, and GRPO. Additionally, we perform experimental studies to investigate key components within this unified framework. In a broad sense, the parameter gradient, denoted by $\theta$, for any given training strategy can be formulated as:

\begin{equation}
    \nabla_{\theta}\mathcal{J}_{\textcolor{red}{\mathcal{A}}}(\theta) = \mathbb{E}[\underbrace{(q,o) \sim \textcolor{red}{\mathcal{D}}}_{Data \ Source}]\left( \frac{1}{|o|} \sum_{t=1}^{|o|}  \underbrace{GC_{{\mathcal{A}}}(q, o, t, \textcolor{red}{\pi_{{rf}}})}_{Gradient \ Coefficient}  \nabla_{\theta}\log \pi_{\theta}(o_t | q, o_{<t})\right).
\label{eq:objective}
\end{equation}

There exist three key components: 1) \textit{Data Source $\mathcal{D}$}, which determines the training data; 2) \textit{Reward Function $\pi_{{rf}}$}, which is the source of the training reward signal; 3) \textit{Algorithm $\mathcal{A}$}: which processes the training data and the reward signal to the gradient coefficient $GC$ that determines the magnitude of the penalty or reinforcement for the data. We analyze several representative methods based on such a unified paradigm:

\begin{itemize}
    \item  \textbf{Supervised Fine-tuning (SFT)}: SFT adapts a pretrained model by additional training on a curated dataset selected by humans.
    \item \textbf{Rejection Sampling Fine-tuning (RFT)}: RFT involves further training of the SFT model on a subset of outputs, which are generated by the SFT model in response to SFT questions and filtered to retain only those with correct answers.
    \item \textbf{Direct Preference Optimization (DPO)}: DPO builds upon the SFT model by conducting fine-tuning with an expanded set of outputs sampled from the SFT model, employing a pairwise DPO loss for optimization.
    \item \textbf{Online Rejection Sampling Fine-tuning (Online RFT)}: In contrast to conventional RFT, Online RFT launches training from the SFT-initialized policy model, iteratively improving it by fine-tuning with new outputs drawn from the live policy model itself.
    \item \textbf{PPO/GRPO}: PPO/GRPO begins with a policy model initialized via SFT, and strengthens the model by leveraging outputs sampled from the evolving policy model throughout the training process.
\end{itemize}

Table \ref{tab:data_policy} provides an overview of the elements comprising these methods. For a comprehensive explanation of the derivation steps, consult Appendix \ref{app:analysis-rl}.

\paragraph{Observation about Data Source} We categorize data sources into two distinct types: online sampling and offline sampling. Online sampling refers to situations where the training dataset is generated from the exploration activities performed by the policy model during its current training phase. In contrast, offline sampling pertains to datasets that originate from the outputs of the preliminary SFT model. Methods such as RFT and DPO utilize the offline approach, whereas Online RFT and GRPO employ the online approach.

As illustrated in Figure \ref{fig:rl-analysis}, our results indicate that Online RFT substantially surpasses RFT across the two evaluated benchmarks. Initially, the performance of Online RFT aligns closely with that of RFT, but as training progresses, Online RFT establishes a marked lead, highlighting the effectiveness of online training methodologies. This trend can be understood intuitively: at the outset, the actor model closely mirrors the SFT model, with sampled data displaying only slight variations. As training advances, more pronounced discrepancies arise in the data sampled from the actor, making the advantages of real-time, online data sampling increasingly apparent.

\paragraph{Observation about Gradient Coefficient} The mechanism by which the algorithm incorporates the reward signal into the gradient coefficient directly impacts model parameter updates. In our study, we differentiate the reward function into two categories: `Rule' and `Model'. The `Rule' approach evaluates the response quality according to answer correctness, whereas the `Model' approach relies on a trained reward model to assign scores to responses. The training samples for this reward model are labeled based on the rule-based judgments. As illustrated by Equations \ref{eq:GC-RFT} and \ref{eq:GC-GRPO}, an important distinction emerges between GRPO and Online RFT: GRPO modulates its gradient coefficient according to the reward scores obtained from the reward model, thereby enabling nuanced positive and negative feedback proportional to the score magnitude. On the other hand, Online RFT does not incorporate this reward-based scaling; it omits penalties for incorrect answers and applies identical positive reinforcement to all responses classified as correct.

As illustrated in Figure \ref{fig:rl-analysis}, GRPO outperforms online RFT, underscoring the effectiveness of modulating positive and negative gradient coefficients. Moreover, the results indicate that GRPO+PS achieves better outcomes than GRPO+OS, demonstrating the advantage of applying fine-grained, step-sensitive gradient adjustments. Additionally, we investigate the impact of iterative RL in our experiments by performing two iteration cycles. Figure \ref{fig:iter-rl} reveals that iterative RL leads to marked performance gains, with the most notable improvement occurring during the initial iteration.

\subsubsection{Why RL Works?} In this study, we apply reinforcement learning utilizing a portion of the instruction tuning data, resulting in a notable improvement over the standard instruction-tuned model. To elucidate the effectiveness of reinforcement learning, we analyze the Pass@K and Maj@K accuracy metrics for both the Instruct model and the RL model across two evaluation benchmarks. As illustrated in Figure \ref{fig:maj-pass}, RL brings about an improvement in Maj@K, while Pass@K remains unaffected. This result suggests that the benefit of RL stems from enhancing the robustness of the model's output distribution—\textbf{specifically, it appears that RL advances performance by increasing the likelihood that a correct answer appears within the TopK choices, rather than by fundamentally strengthening the model’s core abilities.} Consistently, prior work by \citep{wang2023making} highlights a \textbf{misalignment problem} encountered by SFT models in reasoning tasks, demonstrating that their reasoning skills can be improved through various preference alignment techniques \citep{yuan2023rrhf,song2023preference,wang2023making}.

\subsubsection{How to Achieve More Effective RL?}
\label{sec:effec_rl}
Our findings indicate that RL performs effectively on mathematical reasoning tasks. Furthermore, we introduce a comprehensive framework that encapsulates various prominent training strategies, reinterpreting them as direct or approximate applications of RL. According to Equation \ref{eq:objective}, this framework is organized around three essential elements: Data Source, Algorithm, and Reward Function. We also outline several prospective research avenues related to these three aspects.

\paragraph{Data Source} The data source serves as the foundational input for all training approaches. Within the RL setting, we define the data source as unlabeled questions paired with responses generated by sampling from the policy model. Our experiments in this paper limit the data source to the same set of questions utilized during instruction tuning and rely on a straightforward application of nucleus sampling for response generation. We speculate that this limitation may explain why our RL framework yields improvements solely in Maj@K metrics. Future work will investigate the applicability of our RL method to novel, out-of-distribution question prompts, as well as incorporate \textbf{advanced sampling (decoding) strategies}, including tree-search-based methods \citep{yao2023tree}. Furthermore, advancements in \textbf{efficient inference techniques} \citep{xia-etal-2023-speculative,leviathan2023fast,kwon2023efficient,xia2024unlocking}, which shape the exploration capability of policy models, remain highly significant to this endeavor.

\paragraph{Algorithms} The algorithms utilize both data and reward signals to calculate the gradient coefficient, which in turn updates the model parameters. Referring to Equation \ref{eq:objective}, contemporary methods largely place full \textbf{TRUST} in the reward function to adjust the conditional probability—either raising or lowering it—for a particular token. Nonetheless, this reliance on the reward signal is not always justified, as its reliability cannot be guaranteed, especially in the context of highly complex tasks. As an illustrative case, the PRM800K datasets \citep{lightman2023let}—despite extensive annotation by skilled annotators—still exhibit an error rate of about 20\% in the labeling\footnote{\url{https://github.com/openai/prm800k/issues/12\#issuecomment-1728491852}}. Therefore, we intend to investigate reinforcement learning algorithms designed to withstand noise in the reward signals. We are convinced that alignment approaches based on \textbf{WEAK-TO-STRONG} \citep{burns2023weak} principles could introduce significant advancements in learning algorithm design.

\paragraph{Reward Function} The reward function serves as the principal provider of training feedback. In reinforcement learning (RL), it is typically instantiated via a neural reward model. We identify three key avenues for advancing reward model research: 1) \textbf{Improving the generalization capacity of reward models.} It is crucial for reward models to generalize well, enabling them to address out-of-distribution inputs and sophisticated decoding results. Without this, reinforcement learning might only serve to solidify the language models' output distribution, failing to drive real performance enhancements; 2) \textbf{Capturing reward model uncertainty.} Quantifying uncertainty may establish a connection between suboptimal reward models and learning algorithms designed to transition from weak to strong supervision; 3) \textbf{Developing high-fidelity, process-level reward models efficiently}, so as to deliver nuanced and detailed learning signals throughout the reasoning workflow \citep{lightman2023let,wang2023math}.

\section{Conclusion, Limitation, and Future Work}

In this work, we introduce DeepSeekMath, which achieves superior results compared to all other open-source systems on the competitive MATH benchmark, narrowing the performance gap with proprietary models. The training of DeepSeekMath begins with the DeepSeek-Coder-v1.5 7B model and involves ongoing training on a dataset of 500B tokens, with 120B of these being mathematical tokens curated from Common Crawl. Through an in-depth ablation analysis, we demonstrate that web pages are a valuable source of high-quality mathematical content, while data derived from arXiv contributes less than anticipated. We propose Group Relative Policy Optimization (GRPO), an extension of Proximal Policy Optimization (PPO), which enables improved mathematical reasoning skills while reducing memory usage. Experimental results indicate that, even for DeepSeekMath-Instruct 7B, which already performs strongly on evaluation benchmarks, GRPO yields further gains. Furthermore, we offer a unified framework to conceptualize various approaches and outline several promising avenues for future enhancements in reinforcement learning efficiency.

While DeepSeekMath~demonstrates strong results on benchmarks centered around quantitative reasoning, its proficiency in areas such as geometry and theorem proving remains comparatively limited when measured against proprietary models. For example, preliminary evaluations reveal that the model struggles with questions pertaining to triangles and ellipses, which may suggest the presence of biases in the selection of training and fine-tuning data. Moreover, the smaller scale of the model constrains its few-shot learning abilities—DeepSeekMath~does not exhibit notable improvements when provided with few-shot examples, contrasting with GPT-4, which benefits from such inputs. Moving forward, our plan includes refining the data selection process to enhance the overall quality of the pre-training corpus. Additionally, we intend to investigate further directions (Section \ref{sec:effec_rl}) for advancing the reinforcement learning approaches for large language models.

\end{document}